\section[第八章\q 对交易本质的思考]{第八章\\ 对交易本质的思考}

\subsection{交易获利的基石是什么}

{\kaishu
交易者“亏损的天赋”来自于上帝强加给我们的本能，以及社会强加给我们的信念。

\hfill ——金融帝国
}

我们已经探讨了很多关于交易的问题。或许有的朋友还是会有疑问：我所阐述的东西似乎都是建立在经验与感悟下的产物。那么我能否从逻辑上证明我所阐述的东西是正确的呢？或者说，我的交易策略能够获利的基石又是什么呢？

即使是严谨的科学理论，往往也都是从假设开始的。这里我只需要假设，多数人亏损定律。我刚进入股市的时候，就听说过“七赔、两平、一赚”的说法。后来我才知道，这居然是道氏理论三大假设中的说法。很长时间以来，我一直觉得这是一条“不需证实的公理”。但在一次与新手讨论时，他似乎认为交易者整体的战绩远没有如此的悲观。当然，那是在波澜壮阔的大牛市中。或许讨论股票市场时，应该加上“一轮牛熊循环过后”的前提。而对于零和游戏的期货市场来说，多数人亏损定律将会表现得更加明显。

这时一个有意思的问题出现了：无论市场（波动）还是交易者（行为），只要一方是绝对随机而没有偏向的话，那么这条定律就不可能成立！这并不难理解，交易者的行为偏向与市场波动偏向一定是对称的。否则市场波动将无法通过捕捉交易者的行为，而产生“多数人亏损定律”。或者我们可以换一种说法，在那个庄股横行的年代，如果散户的交易行为没有任何的偏向，那么庄家也绝没有可能靠操纵股价走势实现获利。事实上，我根本无法去研究庄家的行为，但我却可以把庄家的行为视为散户集体行为偏向的一种镜像。我们可以将交易者行为中的偏向称为交易者“亏损的天赋”，那么这种天赋又是如何形成的呢？

很显然，交易者获利的本领是需要在市场中，经过长时间的学习才能获得的。而“亏损的天赋”却是在交易者进入市场之前就已经具备的。这种“亏损的天赋”并不是什么特例，而是在所有人身上都普遍存在的。每个人都存在的共性，只可能来自于上帝强加给我们的本能，以及社会强加给我们的信念！而这两点正是交易者“亏损的天赋”的成因。事实上，本书贯穿始终的一条主线就是对交易者心理偏向与思维偏向的讨论。

这并不算是秘密，我没有正统地念过大学，所以数学知识与电脑知识少得可怜。这似乎就注定了，我无法使用复杂的逻辑与烦琐的统计。如果市场的波动是交易者行为的镜像，那么直接研究交易者“亏损的天赋”就会简单得多。请注意，市场既是运动员、又是裁判员。那么交易者整体的任何偏向，无论对错都会成为市场捕捉的目标。比如，如果股民整体都对“业绩是股市永恒的主题”深信不移的话，那么这也可以成为被市场捕捉的一种微观偏向。

为了形象地说明交易者的偏向，我们可以采用“彩票游戏”来说明。如果有一种两位数的彩票，赢家平分输家的下注，并且没有任何收税，那么是否存在某种策略，可以实现稳定盈利呢？中奖的概率是百分之一，而彩票并没有任何概率上的偏向，似乎我们没有办法找到具有正期望收益的游戏策略。事实上，即使游戏本身没有任何偏向，参与者信念上的偏向照样可以让我们找到具有优势的策略。很多人都错误地认为：两期号码相同的概率是万分之一，所以买入上期获奖号码的人会更少。但实际上，本期中奖号码仍旧是任何数字的概率都是百分之一。甚至还会有人可笑地认为，像14（要死）、13（西方人认为不吉利）这样的数字不吉利，所以购买的人同样会很少。那么买入与上期相同的号码、不吉利的号码，虽然不能让你获得任何概率上的优势，但却能够让你在中奖时获得赔率优势。这样一来，在没有税收的情况下，长期使用这种下注策略就能够使你获利。很显然，这就是一种明显的信念偏向。

再举一个例子，如果你每期都买相同号码的彩票，而这期朋友建议你买其他的号吗。那么你会如何选择号码呢？实际上，你的思维很清楚两个不同号码中奖的概率是相同的，但你还是更有可能买入原来的号码。道理很简单，如果你买入新号码，但老号码中奖让你所产生的后悔程度，要大于你买入老号码但新号码中奖时的后悔程度。后悔理论告诉我们，人的行为在很大程度上是为了避免事后后悔的可能性。或许，这就解释了为什么交易者可以接受套牢股票下跌80％的幅度，却接受不了趋势跟踪系统40％的最大潜在资金回撤幅度。因为前者是不作为的亏损，后者是作为的亏损。当作为的亏损发生时，更容易使人感到后悔。美国的一个电视游戏节目也显示出这种倾向：当游戏参与者对三选一的赌博下注以后，主持人排除剩下两个选项中的一个（把不能中奖的选择剔除）时，仍旧有很多人继续坚持开始的选择。很显然，这就是一种明显的本能偏向。

如果我的推论是正确的，那么就解释了：为什么很多交易大师都非常注重心理学和哲学。心理学的作用就是对自己心灵的探索，而哲学又是在锻炼我们用具有效率的思维方式来探索世界。当然，我在本书中对于心理偏向的探讨要比思维偏向多得多。毕竟，探讨心理偏向只需要面对更少的压力。而对于思想来说，除了人的生命以外，最应该值得尊重的就是人的思想。尊重的最高境界，是尊重他人犯错的权利。毕竟，我无法保证自己一定是对的，更不能保证我思想就一定适合你。世界上最自大的行为，莫过于告诉他人应该如何去思考！

\subsection{交易是一件令人快乐的事情吗}

{\kaishu
我经常问自己：自由的迷茫与无奈的安详，究竟哪个更有意义？

\hfill ——金融帝国

当一个人拥有足够的理智和明确的目标后，就会客观地确定唯一最佳路径，从而丧失了随性的自由。

\hfill ——金融帝国
}

记得某位著名的交易者曾经说过：交易是他穿上衣服时，所能够做的最快乐的事情。

有意思的是，我的结论似乎是截然相反的。恐怕这个问题，是论坛上网友们与我争议最大的问题。甚至有人认为，我是在无病呻吟。

按照常理来说，亏损的交易者应该是痛苦的，而获利的交易者应该是快乐的。但我的经历似乎有些同常理不符。当我刚进入股票市场的时候，我总是无法摆脱长期亏损的折磨。但有意思的是，我从来没有放弃过对炒股的那种热情。甚至，当我在人满为患的交易大厅，踮起脚尖看行情的时候，会有一种令人兴奋的激情。在上中专的时候，我一边骑着自行车、一边吃着大饼鸡蛋，就是为了能在中午休息时去股市看上一眼。虽然为了赶时间而狼狈到让面酱流到了裤子上，但我很确信那时是非常快乐的。更有意思的是，当我能够在市场上不断获利，甚至能够让资金曲线平滑地向上移动时，我曾经的那种令人兴奋的激情却消失了！交易对于我来说，几乎成为一种例行公事，而获利的结果并不能使我感到意外。交易对于我来说已经越来越像一份高薪的工作，而这份工作要比通常意义上的工作丢掉更多的自我。更加可怕的是，我追求绝对效率的风格，以及批判性的思维方式被带到了现实生活中。从而在现实生活中更多的去考虑行为的意义，而不是心灵的满足。的确，我似乎已经丧失了自由。当一个人拥有足够的理智和明确的目标后，就会客观地确定唯一最佳路径，从而丧失了随性的自由。甚至我对大众主流的价值观，还有一种本能的恐惧。我似乎在应对真实世界时遇到一种障碍。我已经不习惯依赖任何事物，而心灵上绝对的独立是一种显而易见的痛苦。或许在我听到任何信息的时候，第一反应就是揣摩信息发布者的动机是什么，以及当我相信了这种信息，又会对谁有好处。要知道，将自己的精神交付给别人，就能够得到一种安详的快乐，而我却永远无法得到这种快乐。我经常问自己：自由的迷茫与无奈的安详，究竟哪个更有意义？甚至当我能够不断获利以后，现实生活中娱乐给我带来的快乐较之从前出现了明显的下降。

或许，读者朋友根本无法理解我对交易的感受，但我却可以从心理学的角度来解释这些现象。我相信，人只会“在乎”尚未拥有和即将失去的东西。如果我的观点是正确的，那么一位令你心动的女孩给你带来的那种激情，应该不会延续到婚姻生活当中。人的本能通常都会是先搜寻问题、然后关注问题、最后解决问题，接下来就会去搜索新的问题。而已经得到解决的问题，就会渐渐失去心灵的关注。或许，所谓的激情就是一种对“尚未拥有”的事物的一种关注。如果交易仍旧是一件令你兴奋的事情，那么很可能对于你来说，并没有真正拥有成熟的交易能力。而你真正拥有的东西，应该给你带来的是一种麻木的依赖。我的意思是说，交易获利的能力并不能够给我带来快乐，但失去这种能力就会让我非常痛苦。这很类似于大多数人的婚姻生活：维系婚姻的并非是在一起的快乐，而是分开后的痛苦。

按照认知不协调理论：人们为了减少或避免心理上的不一致，往往当为一件事情付出越多时，他们就越认为这件事有价值，而越有压力为这件事情继续付出。这是一个非常有趣的心理学理论，按照这种理论我们可以得出很多有趣的结论。曾经有位心理学家做过这样的实验：先让一些被试完成一组无聊的实验。当实验结束以后，实验的组织者编出一些理由。希望参与实验的被试告诉后面的人，这是一个令人愉快的实验、非常有趣。为了表示感谢，实验的组织者会给“说谎者”一些金钱的奖励。其中一部分人得到“1美元”，另一部分人得到“20美元”。有意思的是，在事后对这个实验愉悦程度的评价中，“1美元”条件下的被试比“20美元”条件下的被试，评价要高得多。

难道真的存在一种心灵与效率的平衡机制吗？为感情付出而未得到回报的人，往往更加相信付出是值得的，从而付出更多的感情；持有套牢股票而未抛出者，往往更加相信股票的价值，从而更加坚定持股信心；长期买彩票而未中奖者，往往更加相信自己会中奖，从而买入更多的彩票。如果说交易与恋爱是人性最好的两面镜子，那么或许我们可以从恋爱中得到一些启发。往往人对一种行为价值的评价，与付出成正比、与收获成反比。一个女孩越是对你爱搭不理，那么你就越会觉得为了和她在一起付出多少都是值得的。反之，一个女孩越是对你死缠烂打，那么你就越觉得她未必是你最好的选择。

事实上，态度不但可以决定行为，行为同样可以决定态度。没错，态度的改变也可能是由于行为的改变所引起的。为了保持一致性，人们往往不得不去改变他们的信念使之与行为一致。其实道理并不复杂：在我炒股票初期，不断在市场上亏损的时候。我不可能甘心带着亏损离开市场，这一点绝对与大多数人无异。问题是，不断的亏损又是显而易见的事实。为了将亏损赚回来，我必须要相信炒股是可以发大财的。甚至我还能“欺骗”自己的心灵，让自己觉得交易是一件快乐的事情。否则，我还有什么理由继续留在市场中呢？就像上面的心理学实验一样，价格（1美元和20美元）是绝对客观的，而对于价值（愉快程度）却有很大回旋的余地。拿到更多酬劳的人，可以为说谎找到一个很好的理由。而拿到更少酬劳的人，就必须靠调高心灵主观的感受（调低无聊的程度）来与之平衡。

两个可以互为自变量和因变量的东西，就可能形成索罗斯所说的“鞋襻理论”。这很像“煤电联动”现象，发电的成本来自于煤，而煤的开采成本又体现在电上。如果煤涨价，那么就有传递到电价的可能性，而电价的上涨同样可能传递到煤价上面。我见过一些存在赌瘾和彩票瘾的人，他们越是把金钱投入到没有回报的事情上，就会越是坚信只有投入更多才是问题的解决之道。而投入更多的行为，一定会改变（加强）态度，甚至是得到快乐的心理感受。有时候我在想，无论人去做任何没有回报的事情时，心理上一定都会是快乐的。否则，又是什么让他继续做下去呢？或许，当有一天真的得到回报时，才开始反思这种行为是否真的有意义？

我曾经经历过一段让我百思不得其解的感情。上学的时候，有一位对我很好的女孩。她为我付出了很多，甚至还帮我去追我喜欢的女孩。但我却一直都拿她当妹妹，关系从来没有超越友情的范畴。即使是走入社会多年以后，她仍旧对我很好。或许那时谁都能看出她喜欢我，而我也渐渐地喜欢上了她。不知道为什么，我始终没有向她表白的勇气。而她似乎也只是停留在不断暗示的层面，原因或许是她认为女孩子不应该主动表白吧。在这段时间里，我们几乎天天黏在一起，甚至以朋友的身份牵手逛街。后来我终于向她表白，而她也承认一直默默地喜欢我，从来没有真的拿我当过哥哥。最后的结局恐怕超出所有人的想象，我们从确立关系到她选择分手只有短短一周的时间。当时我都快崩溃了，总是不断地逼问她到底是为什么？难道曾经这么多年对我的好，就是为了最后狠狠地耍我一次吗？

有人说恋爱可以使人变成诗人；而失恋可以让人变成哲学家。一段时间以后，我终于理解出现这种结局的合理性。当一个人在没有得到一段感情之前，恐怕是没有资格考虑两个人不适合的地方。唯一感受到的只是对方的好，考虑的就是如何才能在一起。而当位置转变以后，考虑的问题就会立刻变为我们有什么不适合的地方。我觉得在过去很多年当中，一种可怕的正反馈循环在起着至关重要的作用。付出的行为可以增强好感的态度，而好感的态度又会增加付出的行为。当真的有一天，付出有了回报的时候，这种没有根基、靠偏向形成的正反馈环，突然之间彻底崩溃了！有意思吧，这是否非常像索罗斯描述的反射理论呢？或者说，很像市场趋势转折时的真实写照呢？顺便说一句，让我敬佩的人不多，而她的确算是一位。如果没有她当时用理智战胜感情，那么我们最后的结局都要比现在悲惨得多。事实上，我们的差距的确太大了，只要是有眼睛的人都会觉得我们并不适合在一起。她的确很优秀，而那时正巧是我最落魄的时候。我真的感叹“登门槛法”具有的威力，只要有一个开始，然后行为与态度的正反馈就能够使人完全丧失了最基本的理智，甚至对最为明显的差距也视而不见。曾经我一直觉得我的感情经历非常奇特，但后来在身边的朋友中发现越来越多的类似现象。

同样的道理，或许只有能够不断获利的交易者，才有资格考虑交易是否值得的问题。而走上交易之路又会使我们失去些什么。当交易的激情退却之后，生活上遇到的交易综合症又体现了出来。美国有位心理学家曾经做过这样的一个实验：对于彩票中奖者，在生活中看电视、聚会等娱乐活动的心理满足程度进行了调查。结论是，彩票中奖者对于生活中娱乐活动所感受到的快乐程度，明显要低于没有中过奖的人。或许解释的原因就是，生活中的娱乐相对于彩票中奖显得不值一提。当我适应了交易的激情以后，生活中的一切喜怒哀乐都显得不值一提。曾经有位女孩觉得我麻木得像个死人，好像对什么都提不起兴趣来。记得还有一次同一些朋友到游戏城去玩老虎机，赢钱的朋友会表现得兴奋，而输钱的朋友又紧张地希望把钱赚回来。当然，我并没有去玩这种无聊的游戏，只是默默地看着他们。毕竟，在我眼中看到的绝不是输赢的问题，而只是一场具有负期望收益的游戏。相对于参与这种无聊的游戏而言，我则更乐意观察游戏参与者的情绪反应以及他们的下注模式。渐渐地，我开始有些犯困，甚至想找个角落“打个盹儿”。我的朋友向我建议说：“你也玩一会儿吧，无论输赢你都会提起精神的。”而我则告诉他当天在期货市场上亏损的金额，并且解释道：“这点刺激已经无法让我提起神来了。”

前面从得失的角度上阐述了交易对心灵的影响，接下来还是谈谈本人个性化的系统所带来的问题吧。我已经说过，本人的交易策略之所以能够获利，完全是因为对大众偏向的捕捉。事实上，市场很像网络游戏，有人想从中赚钱，而有人想从中取乐。有人天天挂在网上玩网络游戏，只是为了通过枯燥的练级将号码弄得很牛，然后再卖给那些想从网游中取乐的人。

这个世界其实很公平，每个人都会得到一些什么，而又会失去一些什么。至少对于我的交易哲学来说：我可以做一个快乐的交易者，也可以做一个获利的交易者，但我无法将两者兼得。似乎心灵与效率一定是互损的，人做任何事情都会有某些动机：要么是追求效率，要么就是追求心灵。甚至更极端地说，付费让一个人去做他想做的事情，就能够使他丧失了兴趣带来的乐趣。有人说，这个世界上最幸福的事情就是做自己想做的事情，同时又能够赚到钱。而我却觉得，如果把爱好变成了事业，那将会使你失去你的爱好。靠打球取乐的业余爱好者与靠打球为生的专业选手，对待打球这种行为的心理感受很可能是不同的。我相信，靠打球为生的专业选手，很可能小时候同样爱好这种体育运动。区别在于，业余爱好者可以随心所欲地发挥，而职业选手则只能去做有利于增加比分的动作。有意思的是，一位很牛的网络工程师，也和我有相同的感受。当我与他讲述这种感受时，他甚至反问道：“这不是以前我跟你说过的吗？”

我曾经戏称，我所做的工作和高级场所的门童无异。我们都是给予顾客想要的心理满足，同时把他们的钞票带回我的家。唯一不同的地方在于，顾客在得到周到的服务以后，有权决定小费的多少。而对于交易者来说，让我配合着做使他们开心的事情，恐怕昂贵到他们难以想象的地步。难道不是吗？这么多人都想不停地高抛低吸来做价差，从而一次一次得到他们的成就感。那么涨高了大家都卖出，总要有人买吧？很显然，每当我的对手盘体验一次成功的喜悦时，我都要承受一次止损的苦恼。想想都可悲，妓女只不过是出卖肉体，而我却要出卖心灵。

当然，按照边际效率递减规律，长时间做同一件事情，无论是快乐还是痛苦，心灵感觉到的强烈程度都会递减。或许，我早已适应了对于心灵的出卖。记得有一本书说过，当交易者能够将自身心理上的喜怒调整到与大众截然相反时，就会成为一个超人。但我怎么觉得这样的人不像是超人，而更像是怪人呢？

\subsection{交易与生活}

{\kaishu
在知识爆炸的今天，成功的交易（事业）并不取决于你头脑中比别人多些什么，而是取决于你的头脑中能否比别人少些什么。

\hfill ——金融帝国
}

在现实生活中，我认识一位女孩。非常有意思的是，我们很多地方都非常相似，又有更多的地方截然相反。在思想上我们都是喜欢思考的人；而在生活上她又是追求优雅气质的人，而我则是严重的实用主义。她始终觉得我只是一位有特殊追求的混混，或者是愤世嫉俗的怪人。她认为我颠覆了她的一个观点，那就是“腹有诗书气自华”。我总是风趣地回答：我的确看了不少书，但都是批判着看的。

其实，索罗斯在《金融炼金术》中所说的“批判性思维”，只不过是指“逻辑清晰严谨的思考”。甚至有人认为“批判性思维”的字意与原意“逻辑清晰严密的思考”相差太远，应该将其译作“严谨的思考”。但在我看来，思考与批判（质疑）似乎本身就是不可分割的。如果你不去质疑的话，那又该如何去思考呢？学习与思考最大的差异恐怕就体现在批判性，一个是被动的接受，而另一个却是主动的探索。

记得小时候，膳食养生的观点中，主要是讲吃某某东西会有什么样的好处。而现在膳食养生的观点中，往往是讲吃某某东西会对身体产生什么样的坏处。如果说在商品富足的今天，健康的身体并不取决于你多吃些什么，而是取决于你应该少吃些什么的话，那么在知识爆炸的今天，成功的交易（事业）并不取决于你头脑中比别人多些什么，而是取决于你的头脑中能否比别人少些什么。

我的意思是说，你比别人相信更少具有缺陷的观念，才更容易成功！坦率地说，我的交易之路并非是一条不断学习之路，而更像是不断剔除无效观念之路。如果你相信了错误的观念，那么你将会为此付出代价！反过来说，如果你什么都不相信的话，虽然你什么都得不到，但也没有人能够伤害到你。事实上，在我眼中的批判性思维，就是将大众无条件接受的观念，通过逻辑清晰、严谨的思考过滤一遍。

作为交易者，我强烈地建议做上面的功课。但作为一个生活在真实世界的人来说，绝对的理智将会给你带来难以承受的重负。现实与市场有一个相似之处，正确与否并不重要，重要的是顺应潮流。当权证的价格比正股还高时，当豆粕的价格比黄豆还高时，市场的确是错了，但市场却有决定对错的至高无上的权利。恐怕，你所能够做的只是无条件地顺应，即使是顺应一种大众形成的明显错误。倒退几百年，如果你逢人便讲地球是圆的，那么一定会被人当做疯子。

有时在聚会中，我听到谈论的话题无论多么的可笑，仍旧会表现出赞同或者沉默。或许，我可以学得圆滑一些，不去做没有意义的争论。但与众不同的思想一定会体现为与众不同的行为。似乎我在现实生活中，的确像济公说过的那句话一样：别人笑我太疯癫，我笑别人看不穿。很显然，我在交易上养成的思维习惯，已经开始影响到现实生活。而正是这种在市场上可以作为资产的思维方式，在现实生活中却成为了负债。

在现实中生活中，我发现一个有趣的现象：我通常与两种人能够有共同语言，一类是真正把书读明白了的人；另一类则是市井之徒。很显然，前者是靠逻辑来思考问题的答案，甚至另类观点本身就能让他感到兴奋；而后者则是靠自身经历形成的阅历来回答问题，甚至经历本身就能够使我们产生共鸣。遗憾的是，现实中更多受过教育的人则是靠“被灌输的观念”来回答问题。这些人的头脑中被灌输了太多“无需论证的公理”，而从来没有思考过这些所谓的公理是否真的有效。比如说，“业绩好的股票会上涨”似乎也能够成为一种约定俗成的、无需论证的公理。在我刚进入股票市场的时候，同样也是将此视为一种公理，而并没有多问一句为什么。事实上，市场本身就是一种平衡机制。如果绩优股相对于绩差股真的具有更大的上涨潜力，那么必将会吸引更多的资金买入绩优股而抛售绩差股。这样一来市场机制就会通过价格的变化，使它们未来的上涨潜力出现一种平衡。很显然，我成为了这种无效公理的受害者，而多年以后才让我不得不反思这个问题。事实上，现实生活中的人们同样存在更多无效的公理。遗憾的是，市场上约定俗成的公理，我们用简单的逻辑就可以推翻。而在现实中即使是更加明显的无效公理，也很难用清晰的逻辑说服别人，甚至思考这种公理本身就有一种类似于亵渎神灵的压力。很显然，绝对的理智无疑使我体会到一种深层次的孤独。我学不会人云亦云，对于任何观点的接受都必须建立在能够与经历产生共鸣、并且能够通过逻辑来证实有效性的基础之上。最终的结果只能是我与主流人群的思想显得格格不入，而不得不宣称自己做一位隐士。一位朋友马上就开始挖苦我：隐士是那种能够隐藏在人群当中，不会引起人们注意的人。而你穿着的“风格”与疯疯癫癫的习惯真是太引人注目了，这叫隐藏在人群中吗？我突然意识到庄子的那句话：外化而内不化。为了能够捍卫自己独立的思想，恐怕需要依赖于对外在环境的顺应。

对于现在的我来说，似乎已经意识到应该“恶补”在现实中所丢掉的东西。如果说心灵与效率是互损的，那么我既然知道该如何获得效率，就一定会知道如何能够获得快乐。在我看来，一个人如果去做越多“没有意义”的事情，那么这个人就一定会生活的越幸福。

很多人都觉得我不注重礼仪、没有文化修养。我觉得一方面的原因在于，我的谋生手段不需要与任何人接触。外在的礼仪与修养对于我来说，就好像猿猴将没有用的尾巴退化掉一样。虽然“得到别人的好感”对于我来说并非必要的谋生要求，但不修边幅的我却至今都没有找到太太。人的本性真是懒惰的，可以不做的事情就一定不会去做。也许真的是这样，当上天为你打开一扇门，就会为你关上一扇窗。

我想另一个原因就是，我过度追求效率的实用主义。我很喜欢爱因斯坦的一句话，优雅只属于裁缝。从某种意义上来讲，优雅只不过是一种没有效率的病态美。大熊猫为什么比小老鼠优雅呢？难道不是因为大熊猫的无效率吗？优雅的行为注定是无效率的，而可以归纳为礼仪、文化范畴内的行为，注定是“没有意义”的。何为浪漫？在我看来，浪漫不过是出乎常理的行为。从某种意义上来讲，“礼仪与文化”将是“理性与效率”的镜像。或者说，“礼仪与文化”注定是无法用“理性与效率”来解释的行为。比如，过年贴吊钱是一种文化，而这的确是没有效率意义的。那么，为什么这多人都在追求没有效率的东西呢？或许答案就是：没有意义的仪式可以愉悦人类的心灵。这里需要解释一下，我并非是在批判“礼仪与文化”的价值。相反，能够注重过程的优雅，而非追求结果的效率将是一种境界。记得在电视剧中看到霍元甲的一句台词：能够将对手击倒最有效的工具是手枪。或许，交易对于我来说的终极意义将是：在交易上追求“结果的效率”的努力，从而使得我有资格可以在其他全部事情上追求“过程的优雅”。

交易你的计划、计划你的交易，可以让你在市场上获得利润；而生活你的计划、计划你的生活，一定会让你在生活中失去乐趣。其实，我现在已经在很大程度上向大众价值观妥协。虽然我知道“稳定的工作”永远都是虚幻或者廉价的，但我仍旧找了份轻松、稳定的工作。虽然我过去三年在市场上赚到的年平均利润，差不多相当于我在这份工作干到退休的收入（不考虑通货膨胀）。但我仍旧用几倍于放在交易上的精力来应对俗事。相对于从前做生意时应付不完的琐事使我焦头烂额来说，我开始渐渐地喜欢上这种单纯的生活方式。大隐隐于市，应该是一种不错的选择。既然我的事业在虚拟的市场上，那么工作对于我来说就应该是一场有趣的游戏。

顺便说一句，我仍旧过着简朴的生活，或许是因为物价在我眼中以每年两位数的速度上涨。其实，每一位希望自己的事业能够绝对独立的人都会遇到这样的问题，这恐怕就是创业者的悲哀吧。对于打工者而言，无论他们是欠银行百万，还是有百万的存款，下个月的收入都是完全相同的，那么他们的财产就会体现在房子和汽车等资产上面。而对于交易者而言，资金就是事业的血液！如果你多一倍的本金，那么就会增加一倍的利润。每当我的交易进入冬天的时候，我都会不由自主地想到：如果有一天我没有了本金，那么我该靠什么为生呢？

有的时候，我会问自己：我到底是真的喜欢交易呢？还是只不过将交易当做实现生活目标的工具呢？从根本上来说，我是一个非常自我、追求自由的人。但我却不得不心甘情愿地在交易中约束自我。难道追求什么，真的一定会在更高的层次上失去什么吗？我很难说自己是真的喜欢交易，毕竟我从来不会对交易之外的任何事物如此“懦弱”。如果说事业只不过就是工具的话，那么我的确为了交易丢掉了太多太多的生活。至少到现在为止我仍旧是一个失败者，除非我能够学会幸福的生活。接下来，我真的应该随性地去做一些没有意义的事情，寻找一些能够让自己提起兴趣来的爱好。想想看，自己再有一年多就30岁了。如果再学不会享受生活，那么就算十年后事业上能够辉煌又有什么意义呢？

交易与生活的讨论，其实与如何在市场上获利并没有关系。甚至，交易对我生活上的影响有可能与对你的影响截然相反，这取决于每位交易者应对市场的核心哲学。这里我只是想看看大家是否有与我相似的感受。另外就是提醒每一位执迷于交易的朋友，生活才是终极的目的。

\subsection{趋势跟踪系统会失效吗}

{\kaishu
趋势跟踪是一种最不坏的交易策略。

\hfill ——金融帝国
}

这恐怕将是每一位了解趋势跟踪的朋友最为关心的话题了。据说，趋势跟踪的有效性能够追溯到数百年以前，具有价格波动记录的任何市场。对于我来说，人性中的不对称性，无疑将是我们能够找到的最为宏观的规律了。似乎我们有理由认为，趋势跟踪的有效性将是亘古不变的真理。但我并不想给任何人虚假的信念，而更希望站在客观的角度上来分析趋势跟踪的有效性。

对于任何暴利行业来说，只要不存在明显的进入门槛，行业利润都将会向平均利润靠拢。从某种意义上来讲，市场存在的职能就是让暴利行业无法永久地维系。很显然，趋势跟踪系统并不存在任何进入门槛。那么，我们似乎有理由相信，趋势跟踪系统不会永远成为市场上的日不落帝国。

如果说能够维系趋势跟踪有效性的原因来自于亘古不变的人性，那么市场发展的两股潮流或许会在一定程度上削弱有效性的根基。对于中国股市而言，机构投资者掌控资金所占的比例呈现不断增长的态势。很显然，机构投资者相对于个人投资者而言，会具有更小的本能偏向。另一方面，随着我与更多网友进行沟通后发现，不少期货市场上的交易者更加热衷于对交易系统的研究，而研究方式更多地体现在系统的测试上。随着电脑与软件的发展，像“高频交易”这样的模式在十年前简直是无法想象的。听到一些专门玩软件测试的朋友给我讲过的一些奇闻异事，我甚至有一种自己早已落伍的感觉。当然，我对于不了解的事物并不具有发言权，这里就不过多讲述了。但我仍旧感到，如果市场中有大量的人才、投入大量的资金、通过更加先进的电脑，机械而冰冷地寻找市场波动中的偏向，将是一件非常可怕的事情。如果说我的信念来自于多年交易形成的经验，那么电脑狂人对于数据的信念似乎可以跳过我所经历过的必然阶段。或许，他们的努力已经很类似于寻找游戏中Bug（缺陷）的行为，而金融市场只不过是一场大型的网络游戏。

前面已经说过，从对美国股市的研究中不难发现，趋势跟踪的最佳参数在大尺度上不断向迟钝漂移。凭我的经验不难得出这样的结论：1995年以前股票市场的最佳参数应该是10日均线；2000年之前的最佳参数应该是30日均线；从2000年至今的最佳参数应该是60日均线。相对于20世纪90年代初那种完全无噪声的趋势来说，现在趋势跟踪的成本随着趋势噪声的增长而不断增长。无论我们是否愿意相信，趋势跟踪系统所产生的收益都呈现一种衰竭的趋势。我们不难发现，越是成熟的市场，趋势特征就越弱。

很显然，外汇市场中的趋势特征呈现出巨大的噪声特征。或许在周线到月线的级别上，趋势特征仍旧非常明显。但更为复杂的利差问题，以及跟踪趋势所需要备用的庞大资金，使这种大尺度的趋势跟踪的可行性并不大。与此对应的是，只要交易者能够发现任何的微观规律，都将可以使其在短时间内获取暴利。事实上，通过软件测试市场偏向性的主流市场，毫无疑问将是外汇市场。毕竟在100倍的资金杠杆以及市场容量几乎无限的条件下，任何微观偏向都将是价值连城的。

我在很长一段时间内，都是靠申购新股来管理流动资金。随着中国股市走向成熟，已经开始出现越来越多的新股上市首日跌破发行价的现象。即使我没有投入什么精力来研究新股，但也可以轻松地发现一个显而易见的规律：当某日市场有几只新股上市时，开盘价高于发行价一点的股票，其后的走势要相对于开盘价低于发行价一点的股票表现更为强劲。这种现象并不难解释，申购新股者的成本情结似乎就注定了，敢于略微高于发行价开盘的新股将会承接更多的抛盘。一方面获利抛出并不具有任何的压力，另一方面对于开盘后跌破发行价的恐惧又使参与者不得不谨慎对待。

这里大家需要考虑的是：对于这种明显的规律来说，我又能够做些什么呢？或许，我唯一能够做的就是，在中签新股的开盘价低于发行价时赶快卖出。很显然，我根本无法通过这种规律来取利，毕竟股市的规则并不能让我实现强弱对冲。但如果换作期货市场的话，那么我可以在强一些的新股上做多、弱一些的新股上做空。无论市场如何波动，我只不过是在对相对强弱下注。那么我想说明的问题又是什么呢？很简单，相对于“规则限制”的市场来说，寻找“规则开放”市场的波动规律将更有意义。反过来说，越是充分竞争的市场，就越不容易出现明显的偏向。毕竟，任何偏向都会被精明的测试者所快速填补。另一个极端的例子就是：我们谁都知道认沽权证将变为废纸，但我们又能够做些什么呢？或许，这就解释了为什么外汇市场更加不容易形成波动偏向的原因。

从另一个角度来讲，索罗斯的资金量或许只能够在外汇市场施展拳脚，而期货市场的趋势特征似乎对其并没有太大的意义。甚至就连阻击外汇市场的案例来说，也仍旧逃不过三部曲。首先，必须是英国的金融当局犯下错误；然后，必须是金融当局负隅顽抗而知错不改；最后，索罗斯才能够实现阻击。从索罗斯的著作中我们可以发现，当英国政府宣布向外汇市场投入的救市资金时，索罗斯则庆幸这与他准备继续做空的资金相当。我们可以想象：别说英国政府主动改正错误，就算是英国政府不作为，索罗斯也不可能实现成功阻击英镑。道理很简单，做空英镑的资金从全球蜂拥而至，没有英国政府的配合索罗斯根本无法找到对手盘。

前面已经说过，大资金的操盘者并不可能靠试错来实现战略性建仓捕捉到正确的方向。那么从某种意义上来讲，趋势跟踪或许只能够是小资金者的利器。趋势将会是永远存在的，但市场却可以不给予跟踪的机会。如果你看过外汇日内走势中“立杆子”的走势，那么你就会明白这一点。从另一个角度来说，中国股市虽然已经不像成立初期那样，只是纯粹的猫捉老鼠的游戏，但仍旧还无法避免机构投资者捕杀中小投资者的局面。对于期货市场来说，日成交量无疑相对于持仓量高得惊人，这只能说明期货市场中的主力仍旧是中小投资者。我并不否认，中国金融市场会逐渐走向成熟，但这无疑需要很长的时间。就目前来看，中小投资者作为成交量的主流制造者的局面，很难在短时间内改变。

咱们可以假设一下趋势跟踪系统失灵时的市场特征：熊市的成交量超过牛市；下跌以股灾的形式快速实现；换手率大幅萎缩；中小投资者的比例降到很低的水平。如果我们仔细研究美国股市的特征就不难发现：成交量与股指的波动率成正比，而与涨跌几乎没有任何相关性；市场有时会出现由不计代价抛售股票而形成一步到位的暴跌，或者说牛长熊短成为显著的市场特征；中小投资者往往通过购买基金来间接投资股市。即使是这样，美国股市仍旧存在明显的趋势特征。至少我们可以认为，中国股市在没有出现上述特征之前，趋势跟踪系统是不会失效的。

上面我们从宏观的角度上讨论趋势跟踪有效性的问题，那么接下来咱们应该谈论一下微观问题。如果市场陷入长时间的整理阶段，那么趋势跟踪系统将会怎么办？前面已经说过，趋势转折时形成的资金回撤，要比整理行情中的资金回撤幅度大得多。甚至我们还可以通过分散投资，利用品种之间的差异性，使得交易者在“整理的泥潭”中不会出现任何亏损。但仍旧有位朋友向我提出一个很不错的问题：在过去十年中，国内经济快速扩张，而国际市场上也是货币泛滥。或许趋势跟踪策略能够实现庞大的利润，正是因为这些先决条件。问题是，当有一天经济环境发展到平静期，那么长时间的整理是否会使趋势跟踪策略陷入一种可怕的泥潭呢？

毫无疑问，最近多年趋势跟踪策略实现的庞大利润，的确是需要热络的经济环境作为土壤的。或者说，这些庞大的利润在大尺度上并非是一种常态。但请注意，更多的利润完全来自于更多的投入，而不是更少的成本！甚至，随着经济环境的热络，策略亏损期产生的亏损幅度也是在增加的，只不过趋势利润增加得更快罢了。如果我们仔细研究市场整理阶段的波动特征，那么我们不难发现这种特征非常类似于光辐射的特征：能量越高的光辐射，频率就会越高，而波长就会越短。

事实上，大尺度的整理行情，并非是体现在更多的上下震荡次数，而是体现在更慢的震荡频率以及更大的震荡幅度。我们以上证指数为例，自1994年9月1053点至1996年1月513点，长达16个月的整理行情也只不过是一次简单的ABC三波的平坦型整理；自2000年8月2115点至2001年6月2245点，长达10个月的整理行情也不过同样是ABC三波的平坦型整理，以及其后创出新高后就立即见顶的疲弱五浪；而对于1998年和2003年长达两年的整理行情来说，也不过是形成相对复杂的大B浪。请注意，如果将这些整理形态进行等比例缩小，那么就是一种经常出现在日常走势中的形态。很显然，无论是持续多长时间的整理形态，它们都具有本体相似性，从而使它们具有相同的震荡次数。

其实，对于最近多年的热络行情来说，更容易在短时间内出现“恰到好处”的整理行情。而更大尺度的整理行情，反而可以使我们的跟踪尺度出现利润。在最近几年的交易中，我经常遇到短时间内多次左右挨打的尴尬。而在测试2000年的上证指数时却发现，跟随系统只不过会产生几次并不亏损的无效摆动。当然，在更加热络的趋势行情中，整理的尺度会被充分压缩，从而使我们的跟踪尺度可以将其包容。

我的结论是，当经济环境变为更加平静时，趋势利润会快速衰竭，但整理亏损也会出现下降。与此同时，由于趋势利润的衰竭，趋势转折形成的资金回撤也会出现等比例衰竭。如果真的出现这种情况，那么大尺度整理行情所形成的“交易冬天”出现的亏损幅度，并不会高于过去几年“交易冬天”出现的亏损幅度。简单说，亏损周期持续的时间只不过是从两三个月拉长到一年以上，但亏损幅度并不会增加，甚至还可能降低。

事实上，我们能够找到的持续时间最长的整理行情，应该是日经指数自1992年至2000年的整理行情。这与上证指数自2002年至2004年的整理形态非常相像。那么这时又出现一个问题，是什么造就了这种完美的本体相似性呢？或者说，是什么使得江恩的“四次法则”能够在任何尺度下适用呢？我猜测，答案或许就是“学习效应”。

很显然，高抛低吸是一种极具号召力的策略。在2011年的今天，我发现上证指数已经在过去两年多的时间里走出了多次上下整理的行情，所以使用趋势跟踪系统通过股指期货来赌最后的突破。有意思的是，几乎每一次股指的下跌后，都会有朋友向我询问是否可以抄底？虽然我坦诚地说出自己的道理，但更多的人还是忍不住去抄底。甚至他们还信誓旦旦地说，上证指数每一次跌破2700点都会形成底部。虽然我现在还不知晓股市未来会如何运行，但我却可以肯定有更多的人将成为学习效应的牺牲品。我坚信，整理行情在被大家认出之后，一定会意味着整理的完结。

答案很简单，随着股指一次次跌破2700点后被拉起，必将有越来越多的人在下一次跌破2700点时选择抄底。抄底的人数恐怕会随着“跌破、拉起”的次数出现几何倍数的增长。我们不难想象，如果市场出现更多次上下整理的走势，那么必将有更多的人获得更大的利润。毫无疑问，市场的力量绝对不会使这种局面发生。从另一方面来讲，交易者对于“犯错次数”的忍耐力，往往又是相似的。随着一次次在跌破2700点做空而亏损的发生，必将使趋势信徒集体出现心理崩溃。很遗憾，本人在2001年的股市操作中，就成为这种心理崩溃的牺牲品。的确，除非我有强大的信念，否则几乎不可能继续知错不改。

似乎这又能够给我们一些提示。趋势跟踪策略的利润或许在未来会出现衰竭，但绝不可能在宏观上成为一种亏损策略。毕竟，趋势策略的亏损，必定对应着交易者本能的盛宴。这或许就解释了，为什么越长时间的整理行情之后，就越容易出现壮观的趋势行情。悲观地说，趋势跟踪策略最多不过是成为一种“三年不开张、开张吃三年”的策略。恐怕没有哪个现实中的行业，能够接受长时间的不盈利吧？这样一来，就能够将更多的资金挡在趋势跟踪策略的门外。

最后，我们作一个最为悲观的假设：我很确信即使趋势跟踪策略在西方市场失效，但至少在中国市场还要保持上百年的价值。这是由东西方文化差异所决定的，东方文化认为“性本善”；而西方文化则认为“性本恶”。东方人更热衷于劝人向善，但却很少回答该如何应对恶的问题。与此相反，西方人则更加热衷于如何限制恶、消除恶的土壤。这样一来，对于东方人来说就是“趋善避恶”；而对于西方人来说就是“惩恶扬善”。表面上，这与趋势跟踪毫无关系，但文化的不同绝对能够形成不同的思维方式。如果你把“善恶”换成“盈亏”，那么你还会认为这与趋势跟踪无关吗？我相信，东方人一定更热衷于“趋盈避亏”，而不是如何限制亏损。

我们已经讨论了太多关于趋势跟踪策略的问题，不知道大家是否能够发现趋势跟踪策略的精髓呢？趋势跟踪策略只会去回答市场出现什么样的走势才会形成最大的亏损；但却不会去回答市场出现什么样的走势才会形成最大的利润。很显然，趋势跟踪策略是一种“把弊端摆到桌面上”的防守性策略。而其他的策略，恐怕都是考虑如何“把优势摆到桌面上”的问题，同时又不计代价地隐藏弊端。如果你连弊端是什么都不知道的话，那么我很难想象你该如何把游戏进行下去呢？这个世界上不可能有哪种策略是完美的，这里我只能说：趋势跟踪是一种最不坏的策略。很显然，趋势跟踪策略并非是效率最大的策略，但的确是一种弊端最小的策略。

\clearpage
